% Chapter 4b discussion --- this chapter's share of draft \S17 per the cut
% list in Project_specification.md \S5.1: fitting consequences, forward
% connections, and the population-level open problems. The single-cell
% assay predictions and open problems live in Chapter 4a's discussion.

\section{Discussion}
\label{sec:discussion}

This chapter took the completed single-cell theory of Chapters~3 and~4a and
embedded it in a population. The outcome is a renewal model of within-host
infection forced by the intracellular process, exact in its kernels; the
demonstration that the classical constant-rate model is a projection of it,
valid only one exponential phase at a time; an effective-parameter map whose
ends differ by orders of magnitude; the invariance of $R_0$ under the
replacement of budding by bursting; the flooding advantage of bursting at
establishment, and its price in invasion speed; a spectrum of release models
connecting the two ends; and the honest boundary of all of it, drawn against
HIV. What remains is to say what the results imply for the way models are
fitted, how they connect to the chapters that follow, and what is still
open.

\subsection{Consequences for fitting practice}
\label{sec:discussion-fitting}

Three consequences for the way within-host models are fitted to data.
First, $p$ is not a constant of the pathogen: the effective release rate
$p_{\rm eff}(r)$ falls monotonically with the epidemic growth rate, from
$V_\infty/\mean{T_{\rm prod}}$ in the stationary phase towards $\delta$ in
the acute phase, and fits to acute and set-point data should therefore
yield systematically different estimates, in a direction and by an amount
the theory specifies (\S\ref{sec:peff-section}). Second, the eclipse-like
delay usually inserted into within-host models by hand emerges from the
intracellular dynamics with no extra parameter; its duration is the mean
productive lifetime's shadow, $\lambda^{-1}\log\langle X\rangle_{\rm QS}$
(\eqref{prop:lifetime}). Third, the identifiability negative result of
\S\ref{sec:identifiability} says that no fit of $(p,d_{\Icell})$ --- from
any data, at any growth rate --- can recover the three intracellular rates;
a third observable is required, and the theory names the three cheapest
ones: the burst-size dispersion (geometric ratio $1/a$), the no-burst
fraction $b$, and the shape of the release-time density. An estimate of $p$
alone tells one almost nothing about the cell; an estimate of $p$ plus any
one of the three tells one something testable. The single-cell assays that
measure those three observables directly are described in the discussion of
Chapter~4a, and the biological systems where the estimates have already
been pushed to both extremes --- \ft{} and \textit{B.~anthracis} --- are
weighed there as well.

\subsection{Connections forward}
\label{sec:discussion-forward}

The two-type process of Chapter~5 supplies the most immediate continuation.
In the GATE regime $\delta_1=0<\delta_2$ --- no burst before conversion ---
the no-catastrophe probability has a zero initial hazard, then rising:
that is precisely the eclipse-phase signature which within-host models
insert by hand, emerging mechanistically, with the eclipse duration set by
the conversion rate $\nu$. The renewal construction of this chapter is
ready to receive the two-type kernels; what blocks the way is the two-type
calculation itself (below). The budding--bursting comparison of
\S\ref{sec:flooding} finds its evolutionary home in Chapter~6, where the
spectrum of \S\ref{sec:spectrum} --- partial release, reset catastrophe,
self-excitation --- provides the intermediate phenotypes a selection
argument needs; and the nonlinear mechanisms Chapter~7 discusses in the
Komarova vein (immune saturation, local depletion, the ``all run out at
once'' intuition) are exactly the effects the branching-structure
comparison of \S\ref{sec:flooding} deliberately excludes, and should be
named wherever the flooding advantage is invoked.

\subsection{Open problems}
\label{sec:open-problems}

Five questions remain open on the population side.

\begin{enumerate}
\item \textbf{Two-type killed second moments and the two-type $D$.} The
release kernel of a two-type renewal model needs
$\mean{X_a^2\mathbf1\{\tau_c>a\}}$, $\mean{X_aY_a\mathbf1\{\tau_c>a\}}$,
$\mean{Y_a^2\mathbf1\{\tau_c>a\}}$, and the analogue of $D(t)$
(extinction without catastrophe). The moment hierarchy does not close
upward, exactly as in one type; the one-type escape was to differentiate
the Riccati equation, and the two-type route should be to differentiate the
backward system of Chapter~5. This is the main technical obstacle between
this chapter and an eclipse-aware renewal BMVR, and it should not be
assumed easy.
\item \textbf{The flooding boundary against real parameters.} Map
$\delta(\lambda+\mu)>\lambda\mu$ for \yp{} in the macrophage, and compute
the \emph{size} of the advantage from \eqref{eq:flood}: does the pathogen
sit on the advantageous side, and by how much? The estimates discussed in
Chapter~4a --- the near-certain bursting of \ft{}, the clearance-dominated
anthrax case --- suggest that real pathogens sit on both sides of the
boundary, which is what makes the question worth asking.
\item \textbf{Population-level variance.} Propagate $\mathrm{Var}(W)$ of
\eqref{eq:EW2} through the branching structure to a whole-host prediction:
for single-cell assays the first moments suffice, but for whole-host data
the variance is the main stochastic signal the BMVR cannot reproduce. The
integer-valued nature of the offspring count, whose distribution matters
for establishment \cite{williams2024reproduction}, enters through the same
quantity.
\item \textbf{Partial release and the flooding boundary in $\varphi$.}
Where on the boolean continuum of \S\ref{sec:boolean} does the flooding
advantage survive? Does the closed-form structure persist at intermediate
$\varphi$?
\item \textbf{Literature positioning.} Verify whether Carruthers et al.\
\cite{carruthers2020stochastic} already couple an intracellular BDC to a
between-cell model --- the novelty claim of \S\ref{sec:renewal} depends on
the answer --- and position the renewal embedding against the
Pearson--Krapivsky--Perelson \cite{pearson2011stochastic},
Nelson--Gilchrist--Perelson, Gilchrist--Coombs
\cite{gilchrist2006evolution}, and Heffernan--Wahl lines of work, with the
BMVR's attribution settled between McLean \cite{mclean1993balance},
Nowak--Bangham \cite{nowak1996population}, and Perelson et al.
\end{enumerate}

The single-cell open problems --- a conceptual proof that burst size equals
the quasi-stationary distribution, the chained joint laws for general
$\mu$, and the sensitivity of the kernels to logistic intracellular growth
--- are stated in the discussion of Chapter~4a, where they belong. The two
lists together are the research agenda that this part of the thesis leaves
in good order: each problem is named, and each has a route.
